\section{Introduction}

In \emph{The Black Sheep}, Calvino begins from a society in which theft is universal and equality survives only because every household is both predator and prey. Losses and gains offset one another exactly, so the social order remains predatory without becoming stratified. The story changes when one man stays at home. The raid directed at his house fails, the house he would have robbed remains untouched, and a local disturbance in the transfer network appears at once. The later movement toward hierarchy, delegation, and policing follows from that break \citep{calvino_1993}. The opening order can be written as a unit transfer raiding game in which each household chooses in every period whether to stay home or which other house to raid, occupied houses are safe in the benchmark, and concurrent raiders divide a rival loot whose expected share falls with congestion. Under a balanced theft cycle every household leaves home, every unoccupied house is targeted once, and wealth remains level because each household has one incoming loss and one outgoing gain. Universal predation is then compatible with equality through exact accounting, and that opening order survives as a weak Nash equilibrium of the stage game.

In the repeated version, the horizon is infinite, behavior is restricted to stationary Markov strategies, and one household is permanently honest and therefore remains at home in every period. If the remaining households follow the old cycle on the first perturbed date, the predecessor who targets the honest house returns empty handed, the successor escapes the raid that would otherwise have hit him, and the resulting wealth change is local but exact. Once leaving home also exposes current wealth, the successor's windfall can move him above the cutoff at which personal raiding ceases to be privately attractive, so the original balanced cycle no longer survives as a stationary continuation. A richer environment becomes necessary at that point, because exact unit transfers cannot carry hierarchy very far once wealth begins to diverge. The extension therefore preserves the same population, the same public wealth state, and the same logic of explicit target choice while replacing unit loot with proportional loot and enlarging the occupational menu. Rich households compare personal raiding, staying home, hiring thieves, and later hiring guards; poor households compare freelance theft, paid predation, protection work, and withdrawal. Patronage requires delegated predation to reach targets more effectively than freelancers can and poor workers to face outside options weak enough to leave a surplus after wages and contracting cost. Guarding appears later, once concentrated wealth and prey depletion alter the ranking of control technologies. Hierarchy, delegation, and selective protection then emerge as successive environments, each tied to its own payoff comparison and transition condition.

Several literatures meet naturally at that point. Corruption equilibria remain privately coherent because each participant takes as given an order that depends on the participation of others \citep{andvig_moene_1990,bardhan_1997,shleifer_vishny_1993,tirole_1996}. Local perturbations then matter through the same logic that convention and equilibrium selection models use to study movement across sustained orders \citep{young_1993}. Appropriation and defense enter once predation and protection compete for resources under insecure claims \citep{skaperdas_1992,grossman_kim_1995,hirshleifer_1995}. Hierarchy enters once the rich use vertical organization to redirect predatory labor \citep{olsen_torsvik_1998}. Organized enforcement enters once protection stabilizes an unequal order created during the transition itself \citep{olson_1993,north_wallis_weingast_2009,besley_persson_2009,acemoglu_ticchi_vindigni_2011}. The graph of the balanced benchmark also recalls a directed exchange cycle, which is why \citet{shapley_scarf_1974} remains useful as a limited structural reference. The connection stops at the graph. Transfers here are coercive, occupancy changes the opportunity set of others, and rivalry in extraction creates externalities that do not belong to the voluntary exchange environment.

Aggregate wealth cannot be the welfare object because the benchmark preserves it by construction. Period welfare is therefore defined as the sum of concave utilities over wealth minus theft effort costs, contracting costs, guarding costs, and expected punishment losses. That criterion distinguishes pure transfer, inequality, and resource dissipation, and it also keeps the simulations in their proper place. No full Markov perfect equilibrium is solved for the large stochastic state space generated by matching, delegated predation, and protection,\footnote{The theory isolates the margins that can be derived cleanly. The simulation section studies how those margins interact once the extended state space becomes too large for complete analytical characterization.} and production remains outside the model. The simulations accordingly implement the same action logic, targeting rules, and threshold comparisons used in the theory, then measure transition timing, regime prevalence, inequality, and welfare along paths that the analytical results do not characterize in closed form. Section \ref{sec:literature} places the argument in the relevant literatures. Section \ref{sec:model} presents the unit transfer game, the repeated perturbation result, the proportional extension, the contracting problem, the guarding threshold, and the welfare criterion. Section \ref{sec:simulation} states how those analytical objects are mapped into code. Section \ref{sec:results} presents the figures and tables in the order required by the argument. Section \ref{sec:discussion} closes.
